\section{Sumsets: residue intervals, counting and the spectrum}\label{sec:sumsets}

This section proves Theorem~\ref{thm:A}. It turns the intersection form of Theorem~\ref{thm:cm} into an exact count of the undershoot. The masses of $S_1$ and $S_2$ appear as the limits of the totals. A condition on two integer sequences, similar ordering, characterizes equality in the spectral lower bound and is sufficient, though not necessary, for the totals to be monotone: the Cohen--Macaulay set $\{0,2,3,5\}$ has $\Tot(h)=2$ for every $h\ge1$, while its only modulus, $5$, gives the incomparable pairs $(1,1)$ and $(2,0)$.

\subsection{Residue intervals and exact counting}\label{subsec:counting}

A \emph{modulus} for $A_0$ is a positive integer $m$ with $m\mid d$ and $m\in S_1\cap S_2$; $m=d$ is always one. Fix a modulus $m$, let $(w^{(1)}_r)$ and $(w^{(2)}_r)$ be the Ap\'ery sets of $S_1$ and $S_2$ with respect to $m$, and for $0\le s<m$ put
\[
u_s=\frac{s+w^{(1)}_{-s\bmod m}}{m},\qquad v_s=\frac{w^{(2)}_s-s}{m}.
\]
Subscripts of Ap\'ery elements are read modulo $m$. These are integers, $u_0=v_0=0$, $u_s\ge1$ for $s\ge1$ (as $w^{(1)}_{-s}>0$ is congruent to $m-s$), and $v_s\ge0$. For $h\ge0$ put $N=hd/m$.

\begin{lemma}[Residue intervals]\label{lem:residue}
Let $h\ge0$, $m$ a modulus, and $n=km-s$ with $0\le s<m$ and $0\le n\le hd$. Then $n\in I_h$ if and only if $u_s\le k\le N-v_s$. If $A_0$ is CM, the same holds for $hA_0$.
\end{lemma}

\begin{proof}
An integer $x\ge0$ lies in a numerical semigroup containing $m$ if and only if it is at least the Ap\'ery element of its class. The class of $n$ is $-s$, so $n\in S_1$ if and only if $km-s\ge w^{(1)}_{-s}$, that is $k\ge u_s$; and $hd-n=(N-k)m+s\in S_2$ if and only if $(N-k)m+s\ge w^{(2)}_s$, that is $k\le N-v_s$. For CM sets $hA_0=I_h$ (Theorem~\ref{thm:cm}; both are $\{0\}$ for $h=0$).
\end{proof}

For $m=d$ and CM sets this is the coordinate description of the homogenized Ap\'ery set in \cite[Prop.~2.6(3)]{GimenezGonzalezSanchez2023}, read on the line of level $h$. Every multiple $km$, $0\le k\le N$, lies in $I_h$, so every maximal run of integers of $[0,hd]$ outside $I_h$ lies in one \emph{block} $((k-1)m,km)$, $1\le k\le N$. Call $s\in[1,m-1]$ a \emph{hole} of block $k$ if $km-s\notin I_h$. By Lemma~\ref{lem:residue} the holes of block $k$ form
\begin{equation}\label{eq:holes}
H_k=B_k\cup R_{N-k},\qquad B_k=\{s\in[1,m-1]:u_s>k\},\quad R_j=\{s\in[1,m-1]:v_s>j\}.
\end{equation}

\begin{proposition}[Exact counting]\label{prop:counting}
Let $A_0$ be CM, $h\ge1$ and $m$ a modulus. For $1\le k\le N$ and $1\le s\le m-1$:
\begin{enumerate}
\item[(a)] $D_h(km-s)=\#\{t: s\le t\le m-1,\ [s,t]\subseteq H_k\}$;
\item[(b)] $\Tot(h)=\sum_{k=1}^{N}f(H_k)$;
\item[(c)] for integers $0\le T_-\le T_+\le hd$, $\sum_{T=T_-}^{T_+}D_h(T)$ is the number of triples $(k,s,t)$ with $1\le k\le N$, $1\le s\le t\le m-1$, $T_-\le km-s\le T_+$ and $[s,t]\subseteq H_k$.
\end{enumerate}
\end{proposition}

\begin{proof}
(a) If $s\notin H_k$ both sides vanish. Otherwise let $s'$ be the least non-hole of block $k$ above $s$, or $s'=m$. The predecessor of $km-s$ in $hA_0$ is $km-s'$, so $D_h(km-s)=s'-s$, the number of $t\in[s,s'-1]$, which are the $t$ with $[s,t]\subseteq H_k$. (b), (c) Every integer of $[0,hd]$ not divisible by $m$ is $km-s$ for exactly one such pair $(k,s)$, and multiples of $m$ have undershoot $0$; sum (a).
\end{proof}

In (c) only the target $km-s$ is restricted to the window, not $t$: the run below a target may leave the window. Lattice targets are handled by summing (a), targets $T\ge hd$ by Lemma~\ref{lem:elementary}(i) and common divisors by Lemma~\ref{lem:elementary}(ii). The proof uses CM only through $hA_0=I_h$, so the proposition holds for any $A_0$ at a level with $hA_0=I_h$, and for $I_h$ itself at every level.

\subsection{The spectrum and monotone totals}\label{subsec:spectrum}

Put $\beta_k=f(B_k)$ for $k\ge1$ and $\gamma_j=f(R_j)$ for $j\ge0$. Both sequences are nonincreasing and eventually $0$. The \emph{spectrum} of $A_0$ at the modulus $m$ is the multiset of all $\beta_k$ and $\gamma_j$, padded by zeros and written $\lambda_1\ge\lambda_2\ge\cdots$. We call $(u,v)$ \emph{similarly ordered} if $u_s>u_t$ implies $v_s\ge v_t$, that is, if the pairs $(u_s,v_s)$ form a chain in the product order; this is the hypothesis of Chebyshev's sum inequality and the equality case of the rearrangement inequality \cite[Thms.~43, 368 and 369]{HardyLittlewoodPolya1952}. Put $U=\max_su_s$ and $V=\max_sv_s$.

\begin{lemma}\label{lem:spectrum-mass}
$\sum_{k\ge1}\beta_k=\M(S_1)$ and $\sum_{j\ge0}\gamma_j=\M(S_2)$.
\end{lemma}

\begin{proof}
Multiples of $m$ lie in $S_1$ and $S_2$, so $\M(S_i)$ is the sum over blocks $k\ge1$ of $f$ of the gaps of $S_i$ in the block. As in Lemma~\ref{lem:residue}, the gaps of $S_1$ in block $k$ are the $km-s$ with $s\in B_k$, and a position $(k-1)m+s$, $1\le s<m$, lies in $S_2$ if and only if $k-1\ge v_s$, so the gaps of $S_2$ in block $k$ form a translate of $R_{k-1}$.
\end{proof}

\begin{lemma}[Merge]\label{lem:merge}
Let $(x_k)_{k\ge1}$, $(y_j)_{j\ge0}$ be nonincreasing sequences of nonnegative integers, eventually $0$, and $N\ge1$. Then $\sum_{k=1}^{N}\max(x_k,y_{N-k})$ is the sum of the $N$ largest entries of the multiset of all $x_k$ and $y_j$.
\end{lemma}

\begin{proof}
For $z\ge1$ let $p_z=\#\{k:x_k\ge z\}$ and $q_z=\#\{j:y_j\ge z\}$. Among $1\le k\le N$, $x_k\ge z$ for the first $\min(N,p_z)$ indices and $y_{N-k}\ge z$ for the last $\min(N,q_z)$, so $\max(x_k,y_{N-k})\ge z$ for $\min(N,p_z+q_z)$ values of $k$. The multiset has $p_z+q_z$ entries at least $z$. Sum over $z$.
\end{proof}

\begin{theorem}[Spectrum]\label{thm:spectrum}
Let $A_0$ be CM and $m$ a modulus.
\begin{enumerate}
\item For every $h\ge1$, $\Tot(h)\ge\lambda_1+\dots+\lambda_N$.
\item If $(u,v)$ is similarly ordered, equality holds in (1) for every $h$. For $m=d$, equality holds for every $h\ge1$ if and only if $(u,v)$ is similarly ordered.
\item If $(u,v)$ is similarly ordered, then $\Tot(h)=\sum_{\ell=1}^{hd/m}\lambda_\ell$ for every $h\ge0$. Hence $\Tot$ is nondecreasing and discretely concave, $\Tot(h)\le\M(S_1)+\M(S_2)$, and equality holds exactly when $N\ge\max(U-1,0)+V$, which for $m\ge2$ reads $N\ge U+V-1$.
\end{enumerate}
\end{theorem}

\begin{proof}
(1) The functional $f$ is monotone under inclusion, and strictly monotone under proper inclusion, since a point of $Y\setminus X$ is an interval counted by $f(Y)$ only. By Proposition~\ref{prop:counting}(b), $\Tot(h)=\sum_{k=1}^{N}f(B_k\cup R_{N-k})\ge\sum_{k=1}^{N}\max(\beta_k,\gamma_{N-k})$, and Lemma~\ref{lem:merge} identifies the right side.

(2) If $B_k$ and $R_j$ are not nested, there are $s\in B_k\setminus R_j$ and $t\in R_j\setminus B_k$, so $u_s>k\ge u_t$ and $v_s\le j<v_t$; under similar ordering this is impossible, and every block gives equality. Conversely let $m=d$, so $N=h$, and let $u_s>u_t$, $v_s<v_t$. Then $t\ne0$, so $k=u_t\ge1$; put $j=v_s$ and $h=k+j$. Then $s\in B_k\setminus R_j$ and $t\in R_j\setminus B_k$, so $f(B_k\cup R_j)>\max(\beta_k,\gamma_j)$ and (1) is strict at $h$.

(3) By (2) and Lemma~\ref{lem:merge}, $\Tot(h)$ is the sum of the $N=(d/m)h$ largest entries (also at $h=0$). Raising $h$ adds the next $d/m$ entries, which are nonnegative and entrywise at most those added before; with Lemma~\ref{lem:spectrum-mass} this gives the first claims. For $m\ge2$, $B_k\neq\varnothing$ exactly for $1\le k<U$ and $R_j\ne\varnothing$ exactly for $0\le j<V$, so exactly $U-1+V$ entries are positive. For $m=1$ there are no positions and every entry is $0$.
\end{proof}

In words, the bound takes in each block the larger of its two hole sets, one coming from $S_1$ and one from $S_2$, and it is exact in that block when one of the two contains the other. Similar ordering makes this happen in every block, and at $m=d$ nothing else does. For a general modulus the levels at which (1) is strict are determined by the inversions of $(u,v)$.

\begin{proposition}[Strict levels]\label{prop:strict-levels}
Let $A_0$ be CM, $m$ a modulus and $h\ge1$. Then $\Tot(h)>\lambda_1+\dots+\lambda_N$ if and only if
\[
N\in\bigcup_{u_s>u_t,\ v_s<v_t}[\,u_t+v_s,\ u_s+v_t-2\,].
\]
Thus equality holds at every level if and only if no positive multiple of $d/m$ lies in this finite union; for $m=d$ this is similar ordering.
\end{proposition}

\begin{proof}
By the proof of Theorem~\ref{thm:spectrum}, $\Tot(h)$ is strict exactly when some block $k$ has $B_k$, $R_{N-k}$ not nested, that is when there are $s,t$ with $u_t\le k\le u_s-1$ and $v_s\le N-k\le v_t-1$; for a fixed inversion $(s,t)$ such $k$ exists if and only if $N$ lies in the displayed interval (and then $1\le u_t\le k\le N$).
\end{proof}

\begin{example}\label{ex:spectrum}
(a) For $A_0=\{0,2,5,10\}$ and $m=10$ the positive spectrum is $(14,6,2,2)$; at $h=3$ the bound in Theorem~\ref{thm:spectrum}(1) is $22$, while $\Tot(3)=23$. At $m=5$ the pairs are similarly ordered (Lemma~\ref{lem:family-residue} in Section~\ref{sec:family}), so both notions depend on the modulus. (b) The CM set $\{0,2,8,11\}$ has $\M(S_1)+\M(S_2)=18$ and $\Tot(h)=19,19,18,18,\dots$: the totals are not monotone and exceed the mass. There are five inversions $(s,t)$, namely $(6,1)$, $(6,7)$, $(6,5)$, $(4,5)$ and $(10,5)$, with intervals $[1,1]$, $[1,1]$, $[1,2]$, $[2,2]$ and $[2,2]$, whose union is $[1,2]$, so by Proposition~\ref{prop:strict-levels} the inequality of Theorem~\ref{thm:spectrum}(1) is strict exactly at $h=1$ and $h=2$.
\end{example}

\subsection{The threshold}\label{subsec:threshold}

For a CM set and $m=d$ the homogenized Ap\'ery set consists of the points $(u_sd-s,v_sd+s)$, of level $u_s+v_s$ \cite[Prop.~2.6(3)]{GimenezGonzalezSanchez2023}, and the regularity of $\Bbbk[\mathcal S]$ is the largest level \cite[Thm.~3.7]{GimenezGonzalezSanchez2023}:
\begin{equation}\label{eq:reg-cm}
\reg \Bbbk[\mathcal S]=\max_{0\le s<d}(u_s+v_s).
\end{equation}
Under similar ordering one index attains both $U$ and $V$: if $u_s=U$ and $v_t=V$, then either $u_t=U$, or $u_s>u_t$ and $v_s\ge v_t$. So the maximum is $U+V$. (Similar ordering concerns the pairs, not the Ap\'ery points: for $\{0,1,2,5\}$ the indices $s=1,2$ both give $(1,1)$, while their points $(4,6)$ and $(3,7)$ are incomparable.)

\begin{proposition}[Threshold]\label{prop:threshold}
Let $A_0$ be CM and $m\ge2$ a modulus at which $(u,v)$ is similarly ordered. Then $U+V-1=\lceil(c_1+c_2)/m\rceil$, $\Tot(h)=\M(S_1)+\M(S_2)$ if and only if $h\ge\sigma(A)$, and $\reg \Bbbk[\mathcal S]=\sigma(A)+1$. If $m=d$, then
\[
U+V-1=\Bigl\lceil\frac{c_1+c_2}{d}\Bigr\rceil=\sigma(A).
\]
\end{proposition}

\begin{proof}
The Ap\'ery elements with respect to $m$ are $u_sm-s$ for $S_1$ and $v_sm+s$ for $S_2$. Let $s_1$ be the least $s$ with $u_s=U$ and $s_2$ the largest $s$ with $v_s=V$. Since $0\le s<m$, the largest Ap\'ery elements are $Um-s_1$ and $Vm+s_2$, so $c_1+c_2=(U+V-2)m+s_2-s_1+2$ (this also covers $S_i=\N$). As $m\ge2$, $U\ge1$ and $s_1\ge1$; an index attaining both maxima lies in $[s_1,s_2]$; so $2\le s_2-s_1+2\le m$, that is $(U+V-2)m<c_1+c_2\le(U+V-1)m$, and the ceiling is $U+V-1$. As $N=hd/m$ is an integer, $hd\ge c_1+c_2$ if and only if $N\ge U+V-1$. For CM sets $\sigma(A)=\lceil(c_1+c_2)/d\rceil$ \cite[Prop.~3.2]{GimenezGonzalezSanchez2023}, so $h\ge\sigma(A)$ if and only if $hd\ge c_1+c_2$, and Theorem~\ref{thm:spectrum}(3) gives the claim on $\Tot$. For $m=d$ we have $N=h$.

\emph{The regularity.} Let $s$ be an index with $u_s=U$ and $v_s=V$; then $s\ge1$, since $U\ge1$. Put $q=d/m$. Since $m\in S_1\cap S_2$, the Ap\'ery set of $S_i$ with respect to $d$ consists of the elements $w+jm$ with $w\in\Ap(S_i;m)$ and $0\le j<q$: each lies in $S_i$, and an element of $S_i$ congruent to $w+jm$ modulo $d$ is congruent to $w$ modulo $m$, hence equals $w+j'm$ with $j'\ge0$ and $j'\equiv j\pmod q$, so it is at least $w+jm$. Choose $0\le j_1,j_2<q$ with $q-1\le j_1+j_2\le2q-2$ such that $q$ divides $U+V+j_1+j_2$; the range holds $q$ consecutive integers, so exactly one value of $j_1+j_2$ qualifies. Then $x_1=(U+j_1)m-s\in\Ap(S_1;d)$, $x_2=(V+j_2)m+s\in\Ap(S_2;d)$ and $d$ divides $x_1+x_2$, so $x_1$ and $x_2$ are the Ap\'ery elements with respect to $d$ of the classes $-s'$ and $s'$, where $s'=x_2\bmod d$, and the pair of residue $s'$ at the modulus $d$ has sum $(x_1+x_2)/d$. As $U+V+j_1+j_2$ is the least multiple of $q$ that is at least $U+V+q-1$, and $\lceil\lceil x/m\rceil/q\rceil=\lceil x/d\rceil$ for $x\ge0$, \eqref{eq:reg-cm} gives
\[
\reg \Bbbk[\mathcal S]\ \ge\ \frac{U+V+j_1+j_2}{q}=\Bigl\lceil\frac{U+V-1}{q}\Bigr\rceil+1=\Bigl\lceil\frac{c_1+c_2}{d}\Bigr\rceil+1=\sigma(A)+1 .
\]
The reverse inequality is \cite[Thm.~3.16]{GimenezGonzalezSanchez2023}.
\end{proof}

So under similar ordering the total reaches the mass at the level from which $hA_0$ keeps its final shape.

For CM sets with $m=1$ both semigroups are $\N$ and every total vanishes. For $d=1$, $A_0=\{0,1\}$, $\sigma(A)=0$, $\Bbbk[\mathcal S]$ is a polynomial ring of regularity $0$, and both masses and all totals vanish. The nonempty central interval in the definition of $\sigma(A)$ is needed: $\{0,3,4,6\}$ is CM and similarly ordered with $(c_1,c_2)=(6,2)$, and its decomposition with an empty central interval holds already at $h=1<\lceil8/6\rceil$. Similarly ordered CM sets with $d\ge2$ attain the bound $\reg \Bbbk[\mathcal S]\le\sigma(A)+1$ of \cite[Thm.~3.16]{GimenezGonzalezSanchez2023}. The proof uses similar ordering only through an index attaining both $U$ and $V$, so every CM set with such an index at some modulus $m\ge2$ has $U+V-1=\lceil(c_1+c_2)/m\rceil$ and $\reg \Bbbk[\mathcal S]=\sigma(A)+1$. The hypothesis of \cite[Prop.~3.18]{GimenezGonzalezSanchez2023} amounts for CM sets to $s_1=s_2$, which gives such an index. It does not imply similar ordering: $\{0,2,5,10\}$ has $s_1=s_2=7$ at $m=10$ and is not similarly ordered there (Example~\ref{ex:spectrum}).
